\subsection{Long-time mean-field convergence and stabilization}
\label{subsec:long-time}

We now study the effective dynamics as $t\to\infty$.
Finite-dimensional convergence and fixed-horizon trajectory
convergence do not by themselves establish convergence of the
effective process. We first show that exact finite-dimensional
identification times diverge in probability. Under additional
conditions, we then prove convergence of the effective parameters
and scores, together with almost-sure eventual stabilization of
the active pseudo-label state of a representative unlabeled
observation.

Write
$\tilde{\mathbb P}_{\mathcal U}
:=\tilde{\mathbb P}(\,\cdot\mid\tilde\Delta=0)$.

\subsubsection{Nonuniformity of finite-dimensional identification}

Recall the effective active pseudo-label state $\tilde{\mathfrak c}^{\,t}$ from \eqref{eq:mean-field-score}, 
and define the configuration discrepancy by
\[
    \bar h(s,t)
    :=
    \tilde{\mathbb P}
        (\tilde{\mathfrak c}^{\,s}
            \neq\tilde{\mathfrak c}^{\,t} \mid \tilde \Delta = 0).
\]
Thus $\bar h(t,t+1)$ is the probability that a representative
unlabeled observation changes its active pseudo-label state
during the next update.

\begin{proposition}[Divergence of exact identification times]
\label{prop:identification-time-divergence}
Suppose Assumptions~\ref{assumption:asymptotic-regime},
\ref{assumption:quenched-initialization}, and
\ref{assumption:regularity-of-loss-and-penalty} hold.
Then
\[
    \bar h(t,t+1)>0
    \quad \text{for every finite} \quad t\geq1,
\]
and the identification time \eqref{eq:identification-time} diverges in probability:
\[
    T_{\mathrm{id}}
    \xrightarrow{\mathbb P}+\infty
    \qquad\text{as} \qquad d\to\infty.
\]
\end{proposition}

The geometric estimate in Theorem~\ref{thm:st-convergence}
is obtained after exact identification, whose time diverges
in probability with the dimension.
The theorem therefore does not provide a dimension-uniform
onset for this estimate.


\subsubsection{Convergence of the effective process}

\begin{assumption}[Conditions for long-time convergence]
\label{assumption:long-time-loss-penalty}
In addition to
Assumption~\ref{assumption:finite-dim-loss-penalty},
we assume the following:
\begin{enumerate}
    \item There exists a dimension-independent constant
    $B_\ell\in(0,\infty)$ such that
    \[
        |\ell'(y,r)|\leq B_\ell,
        \qquad y\in\{-1,+1\},\quad r\in\mathbb R.
    \]

    \item There exists a dimension-independent constant
    $m_j\in(0,\infty)$ such that
    \[
        j''(u)\geq m_j,
        \qquad u\in\mathbb R.
    \]

    \item The penalty satisfies $j(0)=0$.
\end{enumerate}
\end{assumption}
Assumption~\ref{assumption:finite-dim-loss-penalty}
implies Assumption~\ref{assumption:regularity-of-loss-and-penalty}
with the same constants $L_\ell,L_j$.
Under Assumption~\ref{assumption:asymptotic-regime},
\begin{align}
    p_{\mathcal L}
    :=
    \frac1N\sum_{i\in I_{\mathcal L}}
        \mathbf1_{\{Y_i=+1\}}
    \longrightarrow
    \bar p_{\mathcal L}
    :=
    \tilde{\mathbb P}(\tilde Y = +1 \mid \tilde \Delta = 1).
    \label{eq:limiting-labeled-class-proportion}
\end{align}
For long-time convergence, we require
$\bar p_{\mathcal L}\in(0,1)$.
Introduce the Hilbert space and effective parameter
\[
    \mathcal H
    :=\mathbb R\times L^2(\tilde{\mathbb P}^{[w]}),
    \qquad
    \|(b,u)\|_{\mathcal H}^2
    :=b^2+\tilde{\mathbb E}[u^2],
    \qquad
    \tilde\theta^t:=(\bar b^t,\tilde w^t).
\]
The Hilbert space $\mathcal H$ is distinct from the history
$\sigma$-fields $\mathcal H_t^{[g]}$ and $\mathcal H_t^{[w]}$.
Each effective process is realized jointly over time on its
corresponding probability space from
Definition~\ref{def:mean-field-recursion}.
To specify a dimension-independent step-size condition, define
\begin{align}
    \bar s_\mu
    :=\left(\tilde{\mathbb E}[\tilde\mu^2]\right)^{1/2}
    , \quad
    \bar K_{\mathcal L}
    :=\bar s_\mu+
    \sigma\left(1+\frac1{\sqrt{\bar\rho\bar\delta}}\right)
    , \quad
    \bar K_{\mathcal U}
    :=\bar s_\mu+
    \sigma\left(
        1+\frac1{\sqrt{(1-\bar\rho)\bar\delta}}
    \right), \label{eq:s-mu-k-bar}
\end{align}
and
\begin{align}
    \Lambda_{\mathrm{bd}}
    :=
    \lambda L_j+L_\ell(1+\bar K_{\mathcal L}^2)
    +\pi L_\ell(1+\bar K_{\mathcal U}^2).
    \label{eq:lambda-bd}
\end{align}
We impose the sufficient condition
\begin{align}
    0<\eta<\frac2{\Lambda_{\mathrm{bd}}}.
    \label{eq:mean-field-long-time-step-size}
\end{align}

\begin{theorem}[Convergence of the effective process]
\label{thm:effective-process-convergence}
Assume the hypotheses of Theorem~\ref{theorem:trajectory}
and Assumption~\ref{assumption:long-time-loss-penalty}.
Suppose also that $\lambda>0$, $\pi>0$,
$\kappa_-<0<\kappa_+$, $\bar p_{\mathcal L}\in(0,1)$,
and \eqref{eq:mean-field-long-time-step-size} holds.
Then there exist
\[
    \tilde\theta^\infty
    =
    (\bar b^\infty,\tilde w^\infty)
    \in\mathcal H,
    \qquad
    \tilde r^\infty
    \in L^2(\tilde{\mathbb P}^{[g]})
\]
such that, as $t\to\infty$,
\begin{align}
    \|\tilde\theta^t-\tilde\theta^\infty\|_{\mathcal H}
    \longrightarrow0,
    \qquad
    \|\tilde r^t-\tilde r^\infty\|_
        {L^2(\tilde{\mathbb P}^{[g]})}
    \longrightarrow0.
    \label{eq:effective-score-convergence}
\end{align}
Moreover, there exists a random variable
$\tilde{\mathfrak c}^{\,\infty}\in\{-1,0,+1\}$ such that,
$\tilde{\mathbb P}_{\mathcal U}$-almost surely,
\[
    \tilde{\mathfrak c}^{\,t}
    =
    \tilde{\mathfrak c}^{\,\infty}
    \qquad\text{for all sufficiently large }t,
\]
On labeled observations, choose the extension
\[
    \tilde{\mathfrak c}^{\,\infty}
    :=
    \mathbf1_{\{\tilde r^\infty\geq\kappa_+\}}
    -
    \mathbf1_{\{\tilde r^\infty\leq\kappa_-\}},
    \qquad \tilde\Delta=1.
\]
\end{theorem}

The proof of Theorem~\ref{thm:effective-process-convergence}
proceeds in three steps:
\begin{enumerate}
    \item Passing the finite-dimensional descent inequality to the
    trajectory limit gives a decreasing mean-field potential and
    square-summable parameter increments.

    \item A bound on configuration switching in terms of consecutive
    squared increments gives summable switching probabilities and
    almost-sure eventual stabilization of the representative
    configuration.

    \item Branch-stability estimates combine vanishing increments
    with control of configuration discrepancies to make the effective
    parameters Cauchy. Completeness gives parameter convergence in $\mathcal{H}$,
    and a uniform bound on score differences gives convergence
    of the scores in $L^2(\tilde{\mathbb P}^{[g]})$.
\end{enumerate}
The quantitative estimates and the full proof are given in
Appendix~\ref{sec:proof-long-time}.
The theorem also gives convergence of statistics determined by the
active pseudo-label state.

\begin{corollary}[Stabilization of selection and selected errors]
\label{cor:mean-field-selection-stabilization}
Under the assumptions of
Theorem~\ref{thm:effective-process-convergence},
recall $\bar\omega^t,\bar e^t$ from
Corollary~\ref{cor:mean-field-observables}.
Equivalently,
\[
    \bar\omega^t
    =\tilde{\mathbb P}(\tilde{\mathfrak c}^{\,t}\neq0
        \mid\tilde\Delta=0),
    \qquad
    \bar e^t
    =\tilde{\mathbb P}(\tilde{\mathfrak c}^{\,t}=-\tilde Y
        \mid\tilde\Delta=0).
\]
Then
\[
    \bar\omega^t\longrightarrow
    \tilde{\mathbb P}
        (\tilde{\mathfrak c}^{\,\infty}\neq0 \mid \tilde \Delta = 0 ),
    \qquad
    \bar e^t\longrightarrow
    \tilde{\mathbb P}
        (\tilde{\mathfrak c}^{\,\infty}=-\tilde Y \mid  \tilde \Delta = 0),
\]
and both sequences have finite total variation:
\[
    \sum_{t=1}^{\infty}
        |\bar\omega^{t+1}-\bar\omega^t|
    <\infty,
    \qquad
    \sum_{t=1}^{\infty}
        |\bar e^{t+1}-\bar e^t|
    <\infty.
\]
If, additionally,
\[
    \tilde{\mathbb P}
        (\tilde r^\infty\in\{\kappa_-,\kappa_+\})=0,
\]
then
\[
    \tilde{\mathfrak c}^{\,\infty}
    =
    \mathbf1_{\{\tilde r^\infty\geq\kappa_+\}}
    -
    \mathbf1_{\{\tilde r^\infty\leq\kappa_-\}}
    \qquad
    \tilde{\mathbb P}\text{-almost surely}.
\]
\end{corollary}
